\clearpage\SourceLeaf{12}
\begin{center}\bfseries [NOTIZEN ÜBER CUBISCHE UND BIQUADRATISCHE RESTE.]\end{center}
\begin{center}[1.]\par\smallskip Observatio venustissima inductione facta.\end{center}
\noindent\hspace*{\parindent}2 est Residuum vel non residuum cubicum numeri primi $p$ formae $3n+1$, prout $p$ repraesentabilis est per formam
\[xx+27yy\qquad\text{vel}\qquad4xx+2xy+7yy.\]
\par\noindent\hspace*{\parindent}3 est Residuum vel non residuum, prout $p$ repraesentabilis est per
\[\begin{aligned} &xx+243yy\quad\text{aut}\quad4xx+2xy+61yy\\
\text{vel}\quad&7xx+6xy+36yy\quad\text{aut}\quad9xx+6xy+28yy.\end{aligned}\]
\par\noindent\hspace*{\parindent}5 est $\left.\begin{array}{l}\text{Residuum}\\\text{Nonresiduum}\end{array}\right\}$ si $p$ repraesentatur
\begingroup\small
\[\text{per }\left\{\begin{array}{l}(1,0,675),\ (25,0,27),\ (13,1,52),\ (4,1,169)\\
(7,2,97),\ (9,3,76),\ (19,3,36),\ (25,5,28),\ (25,10,31),\ (27,9,28).\end{array}\right.\]
\endgroup
\begin{center}[2.]\end{center}
\noindent\hspace*{\parindent}Criterium, per quod diiudicatur, utrum numerus $m$ sit residuum cubicum numeri primi $p$ formae $3n+1$.

Fiat
\[4p=aa+27bb\]
